\documentclass[11pt]{article}

\usepackage[margin=1.2in]{geometry}
\usepackage{amsmath,amssymb}
\usepackage{microtype}
\usepackage{booktabs}
\usepackage{xcolor}
\usepackage{mdframed}
\usepackage[colorlinks=true,linkcolor=blue!50!black,urlcolor=blue!50!black]{hyperref}

\newmdenv[backgroundcolor=blue!4,linecolor=blue!40,linewidth=1pt,
  skipabove=8pt,skipbelow=8pt]{keyidea}

\title{Rung One:\\
       \Large A Friendly Guide}
\author{Companion notes to ``A Machine-Verified Exclusion of an
Automatic Collatz Itinerary by Mahler's Method in Degree
One''\thanks{Written in collaboration with Claude (Anthropic). These
notes explain the ideas; the paper and the Lean file contain the
proofs. The Collatz conjecture remains open.}}
\date{August 28, 2026}

\begin{document}
\maketitle

\section{Where we were}

The Collatz game: take a whole number; if it is even, halve it; if it
is odd, triple it, add one, and halve. The conjecture says you always
reach $1$. Nobody can prove it.

Earlier papers in this series set up a \emph{diary} for each orbit:
write $1$ every time the number is odd and $0$ every time it is even.
Two things are known about diaries. First, the diary determines the
starting number completely --- no two starting numbers write the same
diary. Second, if an orbit is going to run off to infinity (the thing
the conjecture forbids), its diary has to contain a lot of ones:
more than about $63\%$ of the entries, forever. Fewer ones means more
halvings, and the number shrinks.

So one way to attack the conjecture is to go through the possible
diaries, from simplest to most complicated, and rule them out one
class at a time. This is the \emph{ladder}. The bottom rung was
already proved: a diary that eventually repeats itself belongs to an
orbit that eventually repeats itself, which is bounded, so it cannot
be a runaway.

The next rung up is the simplest kind of diary that never repeats.
Take the rule
\[
1 \to 1\,1\,0, \qquad 0 \to 0\,1\,1,
\]
start from a single $1$, and apply the rule over and over:
\[
1 \ \to\ 110 \ \to\ 110\,110\,011 \ \to\ 110\,110\,011\,110\,110\,011\,011\,110\,011 \ \to\ \cdots
\]
The limit is an infinite word $w$ that never repeats but is completely
structured: every third letter is a $1$, and every block of three
contains exactly two ones. Two-thirds ones is more than $63\%$, so
nothing we knew forbade a runaway orbit from keeping exactly this
diary. The earlier paper boiled the question down to one number.

\begin{keyidea}
\textbf{The one number.} Think of $w$ as the coefficients of a power
series $F(z) = w_0 + w_1 z + w_2 z^2 + \cdots$. There is a unique
``number'' whose diary is $w$ --- but it is a number in the
\emph{$2$-adic} sense, a kind of infinite binary expansion running off
to the left. It equals $-10 + \tfrac59 F_2(8/9)$, where $F_2(8/9)$
means: plug $z = 8/9$ into the series and add it up $2$-adically
(where $8/9$ is small because $8$ is divisible by a high power of $2$).
An ordinary whole number has a binary expansion that stops. So the
question ``does any whole number have diary $w$?'' is the same as the
question ``is $F_2(8/9)$ a fraction?''
\end{keyidea}

The earlier paper also showed the obvious attack fails. If you cut the
series off after $N$ terms you get a fraction with denominator $9^N$
that agrees with $F_2(8/9)$ to about $3N$ binary digits. To prove a
number is not a fraction by approximation, the agreement has to grow
faster than the denominator, and $2^{3N}$ versus $9^N$ is a losing
race --- the ratio is $0.946$, just under $1$. (And $0.946$ is, once
again, the $63\%$ threshold in disguise.)

\section{What is new}

We prove $F_2(8/9)$ is not a fraction. Therefore no whole number has
diary $w$, and none has a diary that is $w$ with the first few entries
torn off either. No Collatz orbit ever falls into this pattern. This
is the first time a non-repeating diary has been ruled out, and every
step of the proof has been checked by the Lean proof assistant.

The trick that does it is old --- it goes back to Kurt Mahler in 1929
--- and the honest summary is that the earlier paper overestimated
how much of Mahler's machinery would be needed. Here is the idea.

\section{The idea in three moves}

\subsection*{Move 1: a better approximation}

Instead of cutting the series off, find two small polynomials $p_0$
and $p_1$ so that $p_0(z) + p_1(z)F(z)$ starts with a very high power
of $z$. This is a linear-algebra problem: $p_0 = 1 + z^3 - z^4$ and
$p_1 = -1 + z - z^2 + z^3$ have ten coefficients between them, and we
can ask for nine conditions --- the coefficients of $z^0$ through
$z^8$ in $p_0 + p_1 F$ all vanish. They do, and the coefficient of
$z^9$ turns out to be $-1$:
\[
p_0(z) + p_1(z) F(z) = -z^9 + z^{10} + z^{12} + \cdots
\]
This depends only on the first ten letters of $w$, and a computer
checks it in a blink.

Now plug in $z = 8/9$ $2$-adically. Because $8/9$ is $2$-adically tiny
and all the coefficients are whole numbers, the right-hand side is
\emph{exactly} as small as its first term: it is divisible by
$2^{27}$ (that is $8^9$) and not by $2^{28}$. There is no estimating
and no error term; it is exact. This is the point where the $2$-adic
world is kinder than the ordinary one. In Mahler's original setting
one has to work hard to bound this quantity and then work harder to
show it is not accidentally zero. Here both come for free.

\subsection*{Move 2: the functional equation, and the tower}

The rule $1 \to 110$, $0 \to 011$ that builds $w$ translates into an
equation for $F$:
\[
F(z) = \frac{z + z^2}{1 - z^3} + (1 - z^2)\, F(z^3).
\]
It says that $F$ at $z$ is determined by $F$ at $z^3$, and vice versa.
So if $F_2(8/9)$ were a fraction, then $F_2((8/9)^3)$ would be a
fraction, and $F_2((8/9)^9)$, and so on up a tower of points
$(8/9)^{3^k}$ that get $2$-adically smaller and smaller --- and at each
level the fraction gets more complicated. We can keep exact track of
how complicated: at level $k$ the numerator and denominator are at
most about $9^{(5/2)\cdot 3^k}$.

\subsection*{Move 3: the collision}

Now run Move 1 at level $k$ instead of level $0$. The Pad\'e expression
$p_0 + p_1 F$ evaluated at $(8/9)^{3^k}$ is divisible by $2^{27 \cdot
3^k}$ exactly. But if $F_2$ at that point were a fraction $u/D$, the
same expression would be a fraction whose numerator we can bound:
the polynomials contribute at most $9^{4 \cdot 3^k}$ and the tower
contributes at most $9^{3 \cdot 3^k}$ or so, for a total of roughly
$9^{7 \cdot 3^k}$ times a constant. A nonzero whole number divisible
by $2^{27 \cdot 3^k}$ is at least $2^{27 \cdot 3^k}$. So we need
\[
2^{27 \cdot 3^k} \ \le\ (\text{constant}) \cdot 9^{7 \cdot 3^k},
\]
and since $2^{27} = 134{,}217{,}728$ is bigger than $9^7 = 4{,}782{,}969$
by a factor of $28$, the left side wins by a factor of $28^{3^k}$,
which beats any constant once $k$ is large. Contradiction. So
$F_2(8/9)$ is not a fraction.

\begin{keyidea}
\textbf{Why this beats the $0.946$ barrier.} Cutting off the series
costs a factor of $9$ in the denominator for every factor of $8$ in
$2$-adic smallness: ratio $0.946$. The Pad\'e trick costs a factor of
$9^4$ (the degree of the polynomials) for a factor of $8^9$ (the order
of vanishing): ratio $1.89$. Using the functional equation to move up
the tower multiplies both sides by $3$ at each step, so the ratio is
preserved while a fixed overhead is left behind. That squaring of the
approximation quality is exactly what Mahler's method is for; the
earlier paper knew this, but assumed the full transcendence machinery
would be needed. Only the smallest piece of it is.
\end{keyidea}

\section{What was checked, and how}

The Lean file \texttt{Collatz/Rung1.lean} contains no $2$-adic
numbers at all. Every statement above about ``$2$-adically small''
became a statement ``divisible by $2^{\text{something}}$'' about
ordinary integers, and every statement about the series became a
statement about its first $L$ terms, with $L$ arbitrary. The chain is:
\begin{enumerate}
\item The word $w$ is defined by its rule, and its first ten letters
  are computed.
\item The functional equation is proved as an exact identity between
  finite sums (no limits).
\item A whole number with diary $w$ --- or $w$ with $s$ letters torn
  off --- forces a divisibility condition for every prefix length
  (this uses the exact orbit identity from the earlier papers).
\item The divisibility condition climbs the tower.
\item The Pad\'e certificate is checked, and its exact $2$-adic size
  is extracted as a divisibility-plus-oddness statement.
\item The sizes are bounded, and $2^{27} > 28 \cdot 9^7$ does the
  rest.
\end{enumerate}
Lean's kernel accepts the proof using only its three standard
foundational axioms. There is no numerical computation to trust, no
external tool, and no \texttt{sorry}.

\section{What this does and does not say}

It does not say anything about the Collatz conjecture in general.
It rules out one family of diaries --- one infinite word and all its
tails --- out of uncountably many. The value of the result is that it
is the first rung where the diary was genuinely non-repeating, and
that the proof pattern is a template: any word built by a fixed-length
substitution rule in which every block has the same number of ones
gives the same shape of functional equation and the same shape of
certificate, so the same proof should run.

It also does not claim a new theorem in number theory. That $p$-adic
values of Mahler functions at small rational points are irrational is
a known kind of fact; experts would not be surprised by
$F_2(8/9) \notin \mathbb{Q}$. What is new is the connection to
Collatz, the fact that the smallest version of the method suffices,
and the machine-checked proof.

\section{Rung two, begun: the ``don't repeat yourself'' rule}

There is a second result in the paper, of a completely different
flavour, and it needs no power series at all.

\begin{keyidea}
\textbf{The prefix-power rule.} Suppose an orbit's diary starts by
repeating a block of $\ell$ letters over and over --- say $e$ times.
Then the starting number and the number $\ell$ steps later have the
same diary for $(e-1)\ell$ steps, so they agree in their last
$(e-1)\ell$ binary digits. Two different whole numbers that agree in
that many digits differ by at least $2^{(e-1)\ell}$; but $\ell$ steps
of the game can only change a number by a factor of about
$3^{o}/2^{\ell}$, where $o$ counts the odd steps in the block. So
\[
(e-1)\,\ell \;\lesssim\; o \log_2 3 + \log_2 n .
\]
A diary cannot start with a long high power of a short block unless
the starting number is enormous. This is a three-line theorem in Lean.
\end{keyidea}

Now, the threshold: if the block has ones-density $\delta$, the rule
bites when $e > 1 + \delta \log_2 3$. At the $63\%$ line that is
exactly $e = 2$. Above the line one needs a bit more than a square.

Enter the \emph{Sturmian} diaries: the most orderly non-repeating
words there are, the ones you get by writing $1$ each time a point
moving around a circle by a fixed irrational angle crosses a mark.
They have a famous property: they almost repeat with period $q_k$ for
every denominator $q_k$ of the continued fraction of the angle, and the
almost-repetition lasts for about $q_k + q_{k+1}$ letters --- a power
with exponent $2.6$ or more, comfortably above the threshold. Feeding
that into the rule, with careful bookkeeping of how the orbit grows
before the repetition starts, gives the paper's second theorem: no
orbit has a Sturmian diary whose angle has partial quotients $\ge 2$
infinitely often (for angles up to $0.789$; larger partial quotients
handle larger angles). The bookkeeping is machine-checked; the two
facts about circle rotations it relies on are classical
(Lagrange, 1770s; S\'os, 1958) and are stated as hypotheses rather
than proved in Lean. Numerically the rule works for the remaining
``golden'' angles too, up to about $0.9$.

It keeps happening that the $63\%$ line shows up as the exact
threshold of every mechanism we try. That is probably not a
coincidence.

\end{document}
